\clearpage
\renewcommand{\sectioncolor}{quantum}
\section{Five escapes past intuition, with dates}\label{sec:escapes}
\markboth{Five escapes}{}

This is not speculation. Each of the following is a documented case of mathematics deliberately built
past embodied intuition. Each followed \textbf{the same two-step pattern}, and recognising that
pattern is what makes the method repeatable.

\begin{center}
\begin{tikzpicture}[font=\footnotesize]
% ---- timeline axis ----
\draw[slate!55,line width=1.3pt,-{Stealth[length=6pt]}] (0,0)--(16.6,0);
\foreach \x/\y in {0.5/1829,3.6/1868,6.4/1932,8.5/1936,11.3/1990,15.3/2011}{
  \draw[slate!55,line width=0.8pt] (\x,-0.13)--(\x,0.13);
  \node[font=\scriptsize,text=slate,below=3pt] at (\x,-0.1) {\y};}
% ---- event nodes above ----
\node[rounded corners=3pt,draw=quantum,line width=0.9pt,fill=quantum!8,inner sep=5pt,
      text width=3.0cm,align=center,anchor=south] (e1) at (1.0,0.55)
  {\textbf{Non-Euclidean geometry}\\{\scriptsize Lobachevsky 1829\\Bolyai 1832 · Riemann 1854}};
\node[rounded corners=3pt,draw=bio,line width=0.9pt,fill=bio!8,inner sep=5pt,
      text width=2.7cm,align=center,anchor=north] (e2) at (3.6,-0.62)
  {\textbf{Beltrami's model}\\{\scriptsize 1868 --- the \emph{re-grounding}}};
\node[rounded corners=3pt,draw=quantum,line width=0.9pt,fill=quantum!8,inner sep=5pt,
      text width=2.8cm,align=center,anchor=south] (e3) at (6.4,0.55)
  {\textbf{Hilbert space}\\{\scriptsize von Neumann 1932}};
\node[rounded corners=3pt,draw=quantum,line width=0.9pt,fill=quantum!8,inner sep=5pt,
      text width=2.8cm,align=center,anchor=north] (e4) at (9.1,-0.62)
  {\textbf{Quantum logic}\\{\scriptsize Birkhoff \& von Neumann 1936}};
\node[rounded corners=3pt,draw=quantum,line width=0.9pt,fill=quantum!8,inner sep=5pt,
      text width=3.2cm,align=center,anchor=south] (e5) at (12.1,0.55)
  {\textbf{Non-commutative geometry}\\{\scriptsize Connes, 1980s; book 1994}};
\node[rounded corners=3pt,draw=quantum,line width=0.9pt,fill=quantum!8,inner sep=5pt,
      text width=3.1cm,align=center,anchor=north] (e6) at (15.3,-0.62)
  {\textbf{Sheaf-theoretic contextuality}\\{\scriptsize Abramsky \& Brandenburger 2011}};
\foreach \n/\x in {e1/1.0,e3/6.4,e5/12.1}{\draw[slate!50,line width=0.6pt] (\x,0.05)--(\x,0.5);}
\foreach \n/\x in {e2/3.6,e4/9.1,e6/15.3}{\draw[slate!50,line width=0.6pt] (\x,-0.05)--(\x,-0.57);}
% re-grounding arrow
\draw[bio,line width=1pt,-{Stealth[length=5pt]}] (2.68,1.30) to[out=0,in=115] (3.45,0.08);
\node[font=\scriptsize,text=bio,align=center,anchor=west] at (3.62,1.42) {39 years later:\\the re-grounding};
\end{tikzpicture}
\end{center}
\vspace{4pt}
\begin{center}\begin{minipage}{0.93\textwidth}\footnotesize\color{slate}
\textbf{Figure 2.}\; Chronology of the five escapes. Note the gap between the symbolic escape (1829)
and its re-grounding in a drawable model (1868). That gap is the cost of unintuitive mathematics,
measured in decades.
\end{minipage}\end{center}

\subsection{The two-step pattern}

\begin{center}
\begin{tikzpicture}[font=\footnotesize,
  s/.style={rounded corners=4pt,draw=#1,line width=1.2pt,fill=#1!8,inner sep=9pt,
            text width=7.3cm,align=left,minimum height=3.1cm}]
 \node[s=quantum] (a) at (-4.35,0)
   {\textbf{\color{quantum}\normalsize STEP 1 --- Symbolic escape}\\[4pt]
    Negate or delete the assumption that \emph{feels} necessary. Derive consequences. Check for
    contradiction.\\[5pt]
    {\scriptsize You are now computing with something you cannot picture. This is uncomfortable,
    slow, and historically takes decades to be accepted.}};
 \node[s=bio] (b) at (4.35,0)
   {\textbf{\color{bio}\normalsize STEP 2 --- Re-grounding}\\[4pt]
    Build a \emph{new blend} --- a model, a diagram, a calculus --- that makes the structure
    manipulable by hand and eye again.\\[5pt]
    {\scriptsize Beltrami's disk. Feynman's diagrams. The ZX-calculus. This is when the mathematics
    becomes \emph{usable} rather than merely valid.}};
 \draw[-{Stealth[length=7pt,width=6pt]},line width=1.5pt,draw=slate] (-0.6,0)--(0.6,0);
\end{tikzpicture}
\end{center}

\subsection{The cases, one paragraph each}

\begin{auditbox}[title={1. Non-Euclidean geometry \normalfont(Lobachevsky 1829 · Bolyai 1832 · Riemann 1854 · Beltrami 1868)}]
The parallel postulate felt \emph{necessary} precisely because it matched embodied spatial
experience. The escape was symbolic: negate it, derive consequences, find no contradiction. Only
\textbf{39 years later} did Beltrami build the pseudosphere model, and then Klein and Poincaré their
disk models --- new blends that re-grounded the abstract structure in something you could draw.
\end{auditbox}

\begin{auditbox}[title={2. Hilbert space \normalfont(von Neumann, \textit{Mathematische Grundlagen der Quantenmechanik}, 1932)}]
An electron's state is a ray in an infinite-dimensional complex projective space; observables are
self-adjoint operators; measurement is a projection. Not visualisable. But look at the assembly:
\begin{center}
\begin{tikzpicture}[font=\scriptsize,
  q/.style={rounded corners=3pt,draw=#1,line width=0.85pt,fill=#1!9,inner sep=5pt,
            text width=2.85cm,align=center,minimum height=1.15cm}]
 \node[q=slate]   (a) at (0,0)     {vectors as \textbf{arrows}\\{\scriptsize grounded}};
 \node[q=slate]   (b) at (3.35,0)  {inner product as \textbf{angle \& length}\\{\scriptsize grounded}};
 \node[q=sand]    (c) at (6.70,0)  {infinite dimension\\{\scriptsize via the \textbf{BMI}}};
 \node[q=crimson] (d) at (10.05,0) {complexification\\{\scriptsize via the \textbf{rotation metaphor}}};
 \node[q=quantum] (e) at (13.55,0) {\textbf{Hilbert space}};
 \foreach \x/\y in {a/b,b/c,c/d,d/e}{\draw[-{Stealth[length=4pt]},line width=0.85pt,draw=slate] (\x)--(\y);}
\end{tikzpicture}
\end{center}
\vspace{-2pt}
\textbf{The machinery of Part VI is what built quantum mechanics' state space.} The BMI of Chapter 9
and the $90^{\circ}$ rotation metaphor of Case Study 3 are load-bearing components of the formalism
you would use to model an electron in a protein.
\end{auditbox}

\begin{auditbox}[title={3. Quantum logic \normalfont(Birkhoff \& von Neumann, \textit{Annals of Mathematics}, 1936)}]
The lattice of propositions about a quantum system is \textbf{non-distributive}. Distributivity feels
necessary because of the container/collection metaphor: if the marble is in box $A$ or box $B$, it is
in one of them. Delete that, and you have a logic for systems where it is not so.
\end{auditbox}

\begin{auditbox}[title={4. Non-commutative geometry \normalfont(Connes, 1980s; \textit{Noncommutative Geometry}, 1994)}]
\textbf{The sharpest case for your question.} Classical geometry says a space is a set of points, and
functions on it commute because you evaluate at a point. Quantum observables do not commute. Connes'
move: \emph{stop insisting on points.} Take the algebra as primary, drop commutativity, and see what
``geometry'' survives --- you get a geometry with no points, which is what a quantum ``space''
requires.
\begin{center}
\begin{tikzpicture}[font=\scriptsize]
 \node[rounded corners=3pt,fill=crimson!10,draw=crimson,line width=0.9pt,inner sep=6pt,
       text width=6.4cm,align=center] (a) at (-4.0,0)
   {\textbf{the embodied assumption}\\``a space is made of points you could stand at''};
 \node[rounded corners=3pt,fill=bio!10,draw=bio,line width=0.9pt,inner sep=6pt,
       text width=6.4cm,align=center] (b) at (4.0,0)
   {\textbf{delete it, keep the algebra}\\$C^{*}$-algebra, spectral triple, no points};
 \draw[-{Stealth[length=6pt]},line width=1.2pt,draw=slate] (a)--(b)
   node[midway,above=2pt,font=\scriptsize,text=slate]{Move 2};
\end{tikzpicture}
\end{center}
\end{auditbox}

\begin{auditbox}[title={5. Sheaf-theoretic contextuality \normalfont(Abramsky \& Brandenburger, \textit{New J.\ Phys.}, 2011)}]
Quantum contextuality and non-locality turn out to be exactly the obstruction to gluing local
sections of a presheaf into a global section. A piece of mid-century algebraic topology became the
right language for a phenomenon with no classical picture at all --- and, as \S\ref{ss:context}
argued, plausibly the right language for context-dependent molecular identity too.
\end{auditbox}
